\scanpage{12}
8.6. La suite exacte (8.5.9) de systèmes locaux sur $U_d$ correspondant
aux suites exactes (8.5.1) met en évidence que le\deleted{s} systèmes loca\deleted{ux}\inserted{l} \unsure{corres}\inserted{pondant aux}
\deleted{des} $(E(Y)\otimes_W k)_{\mathrm{SS}}$ \deleted{et aux $E(Y,\underline O_Y)_{\mathrm{SS}}$ ont}\inserted{a} en plus d'un sous-système \deleted{resp.
d'un système quotient}\inserted{du premier} de nature essentiellement triviale (car provenant
du corps de base, donc constant du point de vue géométrique) \deleted{une}\inserted{un} \unsure{composant}
\deleted{commun $\operatorname{Coim}\phi=\operatorname{Im}\phi$}\inserted{$E(Y,\underline O_Y)_{\mathrm{ss}}$}, dont on s'attend qu'il donne lieu à une représen-
tation tout à fait non triviale de $G$ (NB si on n'aime pas les groupes
fondamentaux généraux et les systèmes locaux, on peut tout interpréter
en termes de la suite exacte 8.5.1 correspondant au \uline{point générique} $s$ de
l'espace des hypersurfaces de degré $d$ dans $P^r_k$, et remplacer $G$ par le
groupe $\operatorname{Gal}(\overline{k(s)}/k(s))$, -- mais on perd alors les frobenius $\ldots$). On aime-
rait pouvoir préciser ce point. \inserted{(Si $n\geq2$, cette représentation est-elle \uline{irréductible}, \unsure{au moins} \unsure{?})} Par exemple, l'image du groupe \unsure{fondamental}
\inserted{$G$} dans $\operatorname{Aut}(V)$ ($V$ étant le \deleted{\unsure{corps}} vectoriel de dimension finie sur $\underline F_p$ où
s'effectue la représentation envisagée de $G$) est-elle égale à $\operatorname{Aut}(V)$ tout
entier ? \inserted{(Ce n'est pas \deleted{\unsure{exclu} \illegible{1}} lorsque $X=P_k^{\unsure{r=1}}$, cas qu'on peut traiter directement) dès que la condition \unsure{ci-dessus est remplie}}
\deleted{Si oui,} alors dès que \inserted{\unsure{l'on a}} $\dim V\geq2$ (ce qui est sans doute \unsure{vrai}
dès que $d$ est assez grand, $\dim V$ tendant vers $\infty$ avec $d$, cf. 8.5.8),
tout élément de $\underline F_p$ est trace d'un élément de $G$ opérant sur $V$. De ceci,
\inserted{de la ``formule de Woodshole'' \deleted{\unsure{via le|on le} \illegible{4}}}
et du théorème de Čebotareff, on déduirait alors,\deleted{\unsure{,}} si le corps $k$ est fini
qu'on peut, pour n'importe quel ouvert non vide $U$ de l'espace des hyper-
surfaces de degré $d$ dans $P^r_k$ ($d$ étant choisi comme dessus) trouver un
point \deleted{de} $s$ de $U$ à coefficients dans une extension finie $k'$ de $k$, définis-
sant une section $Y$ telle que $\operatorname{card}Y(k')$ soit congru mod $p$ à un élément
de $\underline F_p$ donné à l'avance. Au moment d'écrire ces lignes, et malgré l'assis-
tance de gens compétents comme Katz, il ne semble pas qu'un résultat de
cette nature soit prouvé\deleted{,}\inserted{(} même si $X=P^{2\inserted{=r}}$ i.e. pour des courbes planes \inserted{$Y$} de
grand degré\inserted{)}. \deleted{Ce résultat serait déjà bien intéressant, car il implique-
rait déjà une réponse affirmative à la conjecture du n°6 sur le niveau
de la cohomologie évanescente (dont il est question aussi dans 8.5.8).}
\mspage{11}
